\FloatBarrier
\section{System II: Edwards Aquifer, index well J-17}\label{sec:edwards}

\subsection{Data and specification}\label{sec:edw_data}

The predictand is the calendar-year mean of daily-high elevation at well J-17 (TWDB 6837203, EAA AY-68-37-203), the index well of the San Antonio Pool of the Edwards (Balcones Fault Zone) Aquifer, for 1934--2023 (Table~\ref{tab:edwspec}, Fig.~\ref{fig:edwseries}). Three hydrological objects are read off this one number and should be kept apart: the head is observed at an access point; it indexes a store of water, which is the resource; and the store is replenished and drawn by fluxes, which are the flow. The one-pool water balance closes this loop approximately, in the increment-structure sense used throughout: head change equals weighted fluxes plus a linear drain. It is not a closed mass balance. Annual means use all available daily highs; no year falls below the 240-observation floor (minimum $n=242$ in 1939); missing days are not interpolated; the published pre-continuous composite is used as issued; and the 2023 values carry provisional status. Recharge $R$ is USGS total San Antonio-area recharge ($10^3$ acre-ft\,yr$^{-1}$), estimated by the stream-loss method of \citet{Puente1978} and constructed independently of J-17 \citep{UmphresChoi2025}; pumpage $P$ is well discharge from Edwards Aquifer Authority Table 1 \citep{EAA2024}. Both are constructed fluxes, not observations of J-17, and total-area values mix the San Antonio and Uvalde pools. Spring discharge is stored and not used as a driver. Climate predictors, used only by the climate-informed variants, are September--November Ni\~no~3.4 (HadISST) and calendar-year precipitation in Texas climate divisions 06 and 07 (NCEI nClimDiv). Two thresholds frame the series without entering any forecast: the Edwards Aquifer Authority's Stage~I line (10-day mean J-17 below 660\,ft, a 2007 institutional rule) and the physical level near 618\,ft at which Comal Springs approaches cessation (1956 daily minimum 612.51\,ft). The annual mean is below 660\,ft in 31 of 90 years; the daily minimum is below 618\,ft in one (1956).

\begin{table}[!htb]
\centering\small
\caption{The J-17 specification. D = data, E = empirical construct, M = model, N = normative threshold. The predictand is a measured well series, not a storage reconstruction, a numerical-model inversion, the J-27 Uvalde index or spring discharge.}
\label{tab:edwspec}
\begin{tabularx}{\textwidth}{@{}>{\raggedright\arraybackslash}p{0.24\textwidth}Lc@{}}
\toprule
Field & Contents & Type \\
\midrule
System & Edwards Aquifer, San Antonio Pool, indexed by J-17 & D \\
Predictand $H_t$ & calendar-year mean of daily-high J-17 elevation (ft AMSL), 1934--2023 (90 years) & D \\
Drivers & recharge $R_t$ (USGS, stream-loss method); pumpage $P_t$ (EAA Table 1) & D/E \\
Physical threshold & head at which Comal Springs does not cease ($\approx$618\,ft) & N/E \\
Institutional threshold & EAA Stage I, 10-day mean $<660$\,ft (in force after 2007); the annual-mean indicator is a coarse proxy & N \\
Disturbance & recharge pulses; unmodelled karst; Uvalde--San Antonio exchange & M \\
Service series & USGS 08168710 Comal Springs annual-mean discharge (cfs); scored after retention is frozen & D \\
Fixed windows & drought-of-record drawdown (train 1934--1950, test 1951--1956); drought-of-record recovery (1934--1956 / 1957--1961); pre-permit wet interval (1980--1990 / 1991--1995); critical-period era (1997--2014 / 2015--2023) & --- \\
Rolling origins & minimum 15 training years; $n=75$ ($h=1$), $n=71$ ($h=5$); every model on identical origin sets & --- \\
\bottomrule
\end{tabularx}
\end{table}

\begin{figure}[!htb]
\centering
\includegraphics[width=0.92\textwidth]{edw_fig1_series}
\caption{J-17 annual mean and daily-high range, 1934--2023. 1956 mean 623.15\,ft (daily minimum 612.51\,ft); 1992 mean 691.96\,ft (daily maximum 703.31\,ft); 2023 mean 635.68\,ft. The annual mean is below the 660-ft Stage I line in 31 of 90 years.}
\label{fig:edwseries}
\end{figure}

\subsection{The water-balance ladder}\label{sec:edw_ladder}

\begin{definition}[One-pool water-balance map]\label{def:pool}
With head clipped to $[610,710]$\,ft, the one-pool map is
\begin{equation}\label{eq:pool}
H_{t+1} \;=\; \bigl[\,H_t + \alpha + \beta\tilde R_{t+1} + \gamma\tilde P_{t+1} + \delta H_t\,\bigr]_{\mathrm{clip}},
\end{equation}
where $\tilde R$ and $\tilde P$ are the flux values a forecast substitutes for the unknown year-$(t+1)$ fluxes. Forecasts are issued at the end of year $t$ from $(H,R,P)$ through $t$; $(\alpha,\beta,\gamma,\delta)$ are refitted by least squares on the training transitions at every origin.
\end{definition}

The output clip is $[610,710]$\,ft. It never binds on the observed record (annual means 623--692\,ft); it binds on the 1956-origin persisted-drought M2 path, which is emitted as 617.1, 611.6, 610, 610, 610\,ft for 1957--1961, and at the upper bound on the 1992-origin one-step M2 forecast (raw 719.1, emitted 710\,ft; SI S17). The ladder (Table~\ref{tab:edwladder}) differs from the cod ladder in one structural respect that Section~\ref{sec:synthesis} turns on: the causal rungs must \emph{substitute} something for the target-year fluxes. M2 persists last year's $(R_t,P_t)$; M2m substitutes the training-mean fluxes; M3 adds an AR(1) residual on $\Delta H$ to M2; M4 is a one-year information-delay variant kept as a symmetry control with the cod ladder (J-17 is a telemetered gauge, so the module prices a theoretical lag rather than a constraint of this system). The climate-informed variants replace persisted $R$ by a one-step forecast of $R$ from information known at $t$: an AR(1) on $R$, lagged Ni\~no~3.4, lagged climate-division precipitation, or all three; for $h>1$ the one-step recharge forecast is held constant. The realized-flux oracle receives $(R_{t+h},P_{t+h})$, which is unavailable at the origin; it is an ex-post conditional diagnostic, not a forecast and not eligible for retention. The precipitation oracle is likewise diagnostic. M1 is the univariate AR(1), $H_{t+1}=a+\varphi H_t$, scored as the output-only rung.

\begin{table}[!htb]
\centering\small
\caption{The J-17 ladder. ``Comparator'' is the next-simpler causal model under (E2); M1 has none and is tested against persistence only.}
\label{tab:edwladder}
\begin{tabularx}{\textwidth}{@{}>{\raggedright\arraybackslash}p{0.13\textwidth}>{\raggedright\arraybackslash}p{0.16\textwidth}>{\raggedright\arraybackslash}p{0.2\textwidth}L>{\raggedright\arraybackslash}p{0.14\textwidth}@{}}
\toprule
ID & Class & Fluxes substituted at $t+h$ & Role & Comparator \\
\midrule
persist & baseline & --- & $\hat H_{t+h}=H_t$ & --- \\
mean & baseline & --- & training-window mean & --- \\
M1 & autonomous AR(1) & --- & $H_{t+1}=a+\varphi H_t$ & (persistence only) \\
M2 & causal stock-flow & last $(R_t,P_t)$ persisted & candidate for retention & M1 \\
M2m & climatological stock-flow & training-mean $(\bar R,\bar P)$ & affine AR(1) by class & M1 \\
M3 & residual & as M2 & AR(1) residual on $\Delta H$ & M2 \\
M4 & delay & as M2 & information lag; symmetry control & M3 \\
M2\_Rar, M2\_Renso, M2\_Rprecip, M2\_combo & climate-informed & one-step forecast of $R$ (AR(1) on $R$; Ni\~no~3.4; precipitation; all three) & climate rung & M1 (declared) \\
M2\_oracle & diagnostic & realised target-year $(R,P)$ & ex-post conditional diagnostic; unavailable at origin & --- \\
\bottomrule
\end{tabularx}
\end{table}

\begin{remark}[Class reduction of M2m]\label{rem:class}
Under constant fluxes $\tilde R_{t+1}=\bar R$, $\tilde P_{t+1}=\bar P$, the map \eqref{eq:pool} reduces to $H_{t+1}=(1+\delta)H_t+\mathrm{const}$ and therefore shares M1's forecast function class, though not its estimator: M2m pins its intercept from the in-sample mean fluxes, an additional identifying use of the recharge and pumpage records, while the persistence coefficient $(1+\delta)$ is not mean-flux dependent. M2m passes the H1/H2 point-score gate; the separate class disposition recorded in the broader Aug.\ 27 specification sheet excludes it from the count of additional function classes. That classification is not a failure of (E1) or (E2), and it changes the structural label while leaving the score-gate result intact.
\end{remark}

\FloatBarrier
\subsection{The verdict under Protocol E, and under Protocol C}\label{sec:edw_verdict}

Table~\ref{tab:edwrolling} gives the rolling-origin summary and Fig.~\ref{fig:edwrmse} the per-model RMSE at both horizons. At $h=1$ persistence scores 13.23\,ft and the training mean 16.17\,ft. M1 scores 12.84\,ft and passes (E1); as the simplest output-only rung it has no next-simpler causal comparator, so (E2) is vacuous and M1 is retained by Protocol E's point rule. M2 (14.70), M3 (14.46) and M4 (14.30\,ft) fail (E1): the balance that persists last year's fluxes is worse than persisting the head. M2m (12.28\,ft) passes both H1/H2 inequalities against persistence and M1, then is separately declined as additional structure by the class disposition. The realized-flux diagnostic, at 7.55\,ft, is 43\% below persistence but uses unavailable target-year fluxes. Under Protocol C, restated on these scores, no module is retained: M1's 2.96\% margin is inside the 5\% band and at $h=5$ M1 (21.25\,ft) does not beat persistence (21.11\,ft). The point gate and the class disposition are reported separately because they answer different questions.

\begin{table}[!htb]
\centering\small
\caption{J-17: rolling-origin summary (ft; $n=75$ at $h=1$, $n=71$ at $h=5$, identical origin sets for every row), with Protocol E point-gate status and the separate M2m class disposition. Subscripts give the horizon; $\Delta$ columns are $h=1$ RMSE differences, negative favouring the module. The training mean is the best non-oracle forecaster at $h=5$.}
\label{tab:edwrolling}
\footnotesize\setlength{\tabcolsep}{3.5pt}
\begin{tabularx}{\textwidth}{@{}>{\raggedright\arraybackslash}p{0.19\textwidth}rrrrrL@{}}
\toprule
Model & RMSE$_1$ & MAE$_1$ & RMSE$_5$ & $\Delta$ persist & $\Delta$ M1 & Protocol E decision \\
\midrule
persist & 13.23 & 10.73 & 21.11 & --- & --- & benchmark \\
mean & 16.17 & 13.17 & \textbf{16.80} & $+2.94$ & --- & benchmark \\
M1 AR(1) & 12.84 & 10.72 & 21.25 & $-0.39$ & --- & retained (point rule; margin within noise) \\
M2 persisted fluxes & 14.70 & 11.45 & 33.49 & $+1.47$ & $+1.86$ & rejected: fails (E1) \\
M2m climatological fluxes & 12.28 & 10.22 & 17.44 & $-0.95$ & $-0.56$ & passes H1/H2; separately class-demoted \\
M3 residual & 14.46 & 11.12 & 33.46 & $+1.23$ & $+1.62$ & rejected \\
M4 delay & 14.30 & 11.17 & 33.39 & $+1.07$ & $+1.46$ & rejected \\
M2\_oracle & 7.55 & 5.79 & 10.86 & $-5.68$ & --- & excluded (uses future $R,P$) \\
\bottomrule
\end{tabularx}
\end{table}

\begin{figure}[!htb]
\centering
\includegraphics[width=0.8\textwidth]{edw_fig3_rmse}
\caption{J-17: rolling-origin RMSE by model at $h=1$ and $h=5$. The oracle is a diagnostic ceiling and is excluded from retention. At five years the training mean is the best non-oracle forecaster and M2--M4 are 60\% worse than persistence.}
\label{fig:edwrmse}
\end{figure}

M1's retention is a one-year point-RMSE result and should be read as such. Its 0.39\,ft margin on $n=75$ has a bootstrap interval that covers zero; MAE is a tie (10.72 versus 10.73\,ft), and the ordering reverses at five years. M2m is a second H1/H2 point-gate pass, but its affine-AR(1) class disposition is separate: without that disposition, the literal score gate would count it as retained. The score table therefore does not support the unqualified statement that no stock-flow map passes; it supports the narrower finding that persisting last year's recharge fails on these origins, while the climatological-flux affine map has the best one-step score but is not an additional forecast function class. Neither point ranking is, by itself, a claim of operational value.

\subsection{Uncertainty on the margins}\label{sec:edw_uncertainty}

Table~\ref{tab:edwdm} attaches Diebold--Mariano statistics and moving-block bootstrap intervals to the load-bearing margins, computed from the archived per-origin forecast files; the layer is post-freeze and changes no verdict. Three readings follow. The retention margin itself (M1 $-$ persist, $-0.39$\,ft) is within noise, as are the M2m-over-M1 estimator margin and the M2 causal loss (the last borderline, its interval touching zero); the only one-year margin that separates from zero is M2m over persistence ($-0.95$\,ft, $p=0.003$). The five-year climatology win ($-4.30$\,ft) is robust: its bootstrap interval excludes zero, and the $h=5$ Diebold--Mariano $p$ of 0.10 reflects the known under-sizing of that statistic under overlapping origins. And the climate-rung margins are within noise in both directions. An independent second implementation of the layer (different seed, block length $h$, 20\,000 replicates, different HAC weighting and variance convention; SI S18) reproduces every reading: M1's retention margin has $z=-0.86$ and block interval $[-1.22,+0.56]$\,ft; the M2m edge over persistence has $p=0.001$; the five-year climatology win has $p=0.035$.

\begin{table}[!htb]
\centering\small
\caption{J-17: uncertainty on the load-bearing margins (post-freeze layer; Diebold--Mariano with Newey--West HAC at lag $h-1$; paired moving-block bootstrap, block length 8, 10\,000 replicates, seeded). Rolling origin, $h=1$ unless noted.}
\label{tab:edwdm}
\scriptsize\setlength{\tabcolsep}{3pt}
\begin{tabular}{@{}lrrrcl@{}}
\toprule
Comparison & margin (ft) & DM & $p$ & 95\% block CI (ft) & separates \\
\midrule
M1 $-$ persist & $-0.39$ & $-0.85$ & 0.40 & $[-1.51,\,+0.71]$ & no \\
M2m $-$ persist & $-0.95$ & $-3.07$ & 0.003 & $[-1.45,\,-0.68]$ & \textbf{yes} \\
M2m $-$ M1 & $-0.56$ & $-1.61$ & 0.11 & $[-1.53,\,+0.13]$ & no \\
M2 $-$ persist & $+1.47$ & $+1.27$ & 0.21 & $[-0.02,\,+3.40]$ & no (borderline) \\
mean $-$ persist ($h=5$) & $-4.30$ & $-1.65$ & 0.10 & $[-9.72,\,-2.09]$ & \textbf{yes} \\
M1 $-$ persist ($h=5$) & $+0.15$ & $+0.06$ & 0.95 & $[-5.57,\,+4.94]$ & no \\
M2\_combo $-$ M1 (climate) & $-0.13$ & $-0.21$ & 0.83 & $[-1.37,\,+1.31]$ & no \\
M2\_combo $-$ M2m (nested) & $+0.43$ & $+0.74$ & 0.46 & $[-0.56,\,+1.57]$ & no \\
\bottomrule
\end{tabular}
\end{table}

\FloatBarrier
\subsection{The mechanism: a low-persistence driver, a realized-flux diagnostic, and what climate indices add}\label{sec:edw_mech}

The one-pool increment tracks same-year recharge: $r(\Delta H_t,R_t)=0.74$ on the 1934--2023 record, and the realized-flux diagnostic RMSE is 43\% below persistence at $h=1$ and 49\% below at $h=5$. The same increment is only weakly associated with last year's recharge ($r(\Delta H_t,R_{t-1})=0.06$), while recharge has a low lag-one autocorrelation ($r(R_t,R_{t-1})=0.17$, compared with 0.87 for pumpage). A module that persists $R_t$ therefore substitutes a weakly predictive value and loses to persisting head. Conditioned on the realised target-year fluxes, the same map scores 7.55\,ft; this ex-post conditional diagnostic uses information unavailable at the origin and is not an origin-available forecast. It shows that the fitted map can track the realised increment when supplied its drivers; it does not establish operational skill for predicting the next increment.

The fixed windows make the point visibly (SI S17). The 1957--1961 recovery from the drought of record is a recharge pulse ($R_{1956}=43.7$, $R_{1957}=1142.6\times10^3$ acre-ft): persistence stays at the 1956 floor (RMSE 43.6\,ft), M2 persists drought recharge and falls further, into the clip (55.3\,ft), and the oracle tracks the rise (12.3\,ft). The 1992 peak and the 2015--2023 window repeat the split. On the 1951--1956 drawdown M2 ``beats'' the oracle (18.1 against 19.7\,ft), but for a reason that undermines it: the map trained on 1934--1950 has the wrong sign on pumpage ($\gamma=+0.021$), because pumping rose as the drought deepened and least squares aliases the human response into the state transition.

The climate-informed variants ask whether signals observable at the annual origin (September--November Ni\~no~3.4, lagged climate-division precipitation, and an AR(1) on $R$) supply a useful recharge forecast (Table~\ref{tab:edwclimate}). On the recharge target, the climate modules score 528--545 against the $556\times10^3$ acre-ft climatology; persisted recharge scores 702. Under the pre-score Pass~2 H1/H2 point gate, which compares against persistence and M1, Renso, Rprecip and combo pass at $h=1$; Rar fails. These point-gate outcomes are distinct from the archived class assessment, which demotes the three passes as M2m-plus-intercept variants rather than additional head-forecast structure, leaving \texttt{retained\_as\_structure} empty. Against M2m as an alternative comparator none clears the strict point gate; this alters the point-gate count but not the separate class result. The edges over M1 are only 0.02--0.13\,ft; the reported combined-rung uncertainty interval covers zero (Table~\ref{tab:edwdm}). All four climate variants are 3--6\,ft worse than persistence at $h=5$, in part because the one-step recharge forecast is held constant over the horizon. The Pass~2 score outcomes are reported as written, while the uncertainty and longer-horizon results qualify interpretation rather than retroactively changing the gate. Fixed-window results, including the one window on which two variants edge M1 by about a foot, are in SI S22. Same-year precipitation, unavailable at the origin, yields 10.56\,ft as a conditional diagnostic; it remains above the realized-flux diagnostic because the linear rain map misses 1957-scale extremes. Neither diagnostic forecasts next year's recharge.

\begin{table}[!htb]
\centering\scriptsize\setlength{\tabcolsep}{2pt}
\caption{J-17: climate-informed recharge, rolling origin on the same sets as Table~\ref{tab:edwrolling}. Head RMSE in ft; recharge RMSE in $10^3$ acre-ft. The final column separates point-gate status from the archived structural-class disposition; neither is an uncertainty test. Renso, Rprecip and combo pass the Pass~2 point gate but are class-demoted; M2m's separate Pass~1 result is shown for comparison.}
\label{tab:edwclimate}
\begin{tabularx}{\textwidth}{@{}>{\raggedright\arraybackslash}p{0.20\textwidth}rrrr>{\raggedright\arraybackslash}X@{}}
\toprule
Model & $H$, $h{=}1$ & margin vs.\ M1 & $H$, $h{=}5$ & $R$, $h{=}1$ & Point gate / class status \\
\midrule
persist $H$ / persist $R$ & 13.23 & $+0.39$ & \textbf{21.11} & 702 & benchmark \\
M1 & 12.84 & --- & 21.25 & --- & Pass~1 retained rung; Pass~2 comparator \\
M2\_Rar & 13.25 & $+0.41$ & 25.38 & 561 & fails Pass~2 H1/H2 \\
M2\_Renso & 12.82 & $-0.02$ & 24.42 & 528 & passes Pass~2; class-demoted \\
M2\_Rprecip & 12.80 & $-0.04$ & 25.38 & 545 & passes Pass~2; class-demoted \\
M2\_combo & 12.71 & $-0.13$ & 26.88 & 538 & passes Pass~2; class-demoted \\
M2m climatological fluxes & 12.28 & $-0.56$ & 17.44 & --- & passes Pass~1; separately class-demoted \\
recharge climatology & --- & --- & --- & 556 & benchmark \\
precipitation oracle (year $t{+}h$) & 10.56 & --- & 16.91 & \textbf{354} & conditional diagnostic; future input \\
\bottomrule
\end{tabularx}
\end{table}

\subsection{Post-freeze layers: the service series and the pumpage counterfactuals}\label{sec:edw_post}

Two layers were scored after retention was frozen and change no verdict; both are reported in full in the SI. The Comal Springs discharge series is a near-linear channel of the same state ($r=0.986$ with J-17): rescoring the ladder on the same origins gives a broadly similar ordering (persist 71.9, M1 69.0, M2m 68.7, M2 74.8\,cfs; realized-flux diagnostic 45.3), as expected from a near-linear channel but not mathematically guaranteed. The channel is least redundant with head in the drought tail, where the fitted rating curve predicts about 97\,cfs at the 1956 mean against an observed 32 and fails to signal cessation (SI S20). The pumpage counterfactuals replay actual recharge through the short pre-permit map from the 1990 head; they are illustrative simulations, not forecasts or policy estimates. Across the four paths, the 2023 head spans 630.9--646.8\,ft, comparable to the map's 7.2--14.2\,ft RMSE against observations. This makes pumpage a secondary lever within this fitted map; the coefficient is subject to the drawdown-window simultaneity caveat and no scenario repairs the ladder's primary failure (SI S21).

\subsection{Reading the aquifer verdict}\label{sec:edwreading}

Three statements survive the uncertainty layer. First, the causal balance that persists last year's recharge is not a one-year forecast winner on this origin set; the evidence points to low recharge persistence rather than an inability of the map to track a realised transition. The 7.55\,ft realized-flux score is conditional on target-year drivers, so it does not measure prospective forecast skill. Second, M1's Protocol E retention is a weak-autocorrelation result, not evidence that output-only models are generally preferable. An estimated AR(1) nests persistence at $\varphi=1$; as the head series approaches that limit the forecasts converge and the AR(1)'s advantage shrinks toward estimation noise. Here the head is persistent but far from a unit root (lag-one autocorrelation 0.63; fitted $\hat\varphi=0.66$), leaving a real but slim point margin.  The fitted $\hat\varphi=0.66$ also has a physical reading: under first-order Darcy proportionality, spring discharge above the spring level implies an affine autonomous form $H_{t+1}\approx(1-k)H_t+kH_s$, with $\hat k\approx0.34$\,yr$^{-1}$. Third, on these rolling origins the training mean beats persistence at five years (16.80 versus 21.11\,ft; bootstrap interval excludes zero). This is a benchmark-specific result for a rapidly recharged, institutionally bounded aquifer, not a general multi-year rule; it need not transfer to a non-stationary fossil aquifer under sustained depletion. Section~\ref{sec:synthesis} states the narrower conditions under which a mean-reverting baseline is informative.
